\documentclass[12pt]{article}
\usepackage[utf8]{inputenc}
\usepackage[T1]{fontenc}
\usepackage[margin=1in]{geometry}
\usepackage{setspace}
\usepackage{xcolor}
\usepackage[colorlinks=true, urlcolor=blue!60!black, linkcolor=blue!60!black]{hyperref}
\setlength{\parindent}{0em}
\setlength{\parskip}{0.6em}
\renewcommand{\baselinestretch}{1.0}
\begin{document}
\begin{center}
{\LARGE \textit{Trump v.\ Hawaii}, 585 U.S.\ 667 (2018)\par}
\vspace{0.5em}
{\large Case Brief --- Two-Page Edition\par}
\end{center}

\section{Memory Jogger}
The Court upheld the President's entry restrictions on nationals of mostly Muslim-majority countries under \S~1182(f) and rational-basis review, and repudiated \textit{Korematsu} (\href{https://supreme.justia.com/cases/federal/us/585/17-965/}{Justia}).

\section{Detailed Case Facts}
After a worldwide review of foreign vetting practices, Proclamation No.~9645 restricted entry from eight countries, with exemptions and waivers. Challengers cited the President's call for a ``total and complete shutdown of Muslims entering the United States.''

\section{Procedural History}
The district court entered a nationwide preliminary injunction; the Ninth Circuit affirmed on statutory grounds.

\section{Judicial Votes}
Reversed, 5--4. Roberts wrote for Kennedy, Thomas, Alito, and Gorsuch; Kennedy and Thomas concurred; Breyer (with Kagan) and Sotomayor (with Ginsburg) dissented.

\section{Holding}
Section 1182(f) authorized the Proclamation, and under \textit{Mandel} and rational-basis review it plausibly rested on national security rather than religious animus.

\section{Analysis of Opinions}
\textbf{Majority:} The statute's only condition, a finding of detriment, was met; \S~1152 governs visas, not entry. The Proclamation was neutral on its face, rested on an interagency review, and exempted most Muslims, so it could be understood apart from animus. \textbf{Kennedy:} officials remain bound by the Constitution. \textbf{Thomas:} universal injunctions are suspect. \textbf{Breyer:} the waivers may be ``window dressing.'' \textbf{Sotomayor:} a reasonable observer would see religious hostility, and the Court ``redeploys the same dangerous logic'' as \textit{Korematsu}.

\section{Examples: Future Applications}
\textbf{Same-Side:} a vetting-based suspension of tourist visas despite hostile campaign statements. \textbf{Opposite-Side:} a proclamation barring adherents of a named religion. \textbf{Fence-Sitter:} waivers that exist only on paper.

\section{Critique}
Critics contrast \textit{Masterpiece Cakeshop}'s treatment of hostile remarks and question repudiating \textit{Korematsu} while deferring at the border.

\section{Key Quotations}
\begin{quote}
``The forcible relocation of U.S. citizens to concentration camps, solely and explicitly on the basis of race, is objectively unlawful and outside the scope of Presidential authority. ... \textit{Korematsu} was gravely wrong the day it was decided, has been overruled in the court of history, and---to be clear---`has no place in law under the Constitution.'\,''
\end{quote}
\noindent --- Chief Justice Roberts, for the Court, 585 U.S.\ at 710.

\end{document}
